\documentclass[conference]{IEEEtran}
\usepackage[T1]{fontenc}
\usepackage{amsmath,amssymb,amsthm}
\usepackage{booktabs,tabularx,array}
\usepackage{xurl}
\usepackage[hidelinks]{hyperref}
\hypersetup{pdftitle={Solidity Semantics for Compiler Verification},pdfauthor={OKcontract}}
\Urlmuskip=0mu plus 2mu\relax
\newcommand{\W}{\mathbb W}
\newcommand{\N}{\mathbb N}
\newcommand{\eval}[1]{\mathopen{[\![}#1\mathclose{]\!]}}
\newcommand{\src}[1]{\path{#1}}
\newcommand{\code}[1]{\texttt{#1}}
\newtheorem{proposition}{Proposition}
\title{Solidity Semantics for Compiler Verification:\\Effect Orders, Physical Aliases, and Observations}
\author{\IEEEauthorblockN{OKcontract Labs}
\IEEEauthorblockA{Lean formalization for Solidity 0.8.36 and Cancun\\12 November 2026}}
\begin{document}
\maketitle
\begin{abstract}
Compiler verification requires source semantics that retain observable effect-order
choices, aliasing, and rollback while keeping proof and resource boundaries
explicit. We present Solidity Spec, an executable Lean development targeting
Solidity 0.8.36, its IR compiler pipeline, and Cancun. Its resolved continuation machine selects
eager child expressions independently, operates over a shared physical aggregate
heap, and distinguishes source failure from unsupported behavior and host exhaustion.
The architecture separates checked entry points, fixed implementation policies,
and compatibility imports.
The release combines explicit theorem contracts, admission counterexamples,
independent compiler and EVM checks, and quantified source/target pilots, including
an order-sensitive success-versus-rollback result. External execution checks
assess agreement with Solidity implementations, while configuration invariants
establish admitted-state safety for a restricted profile. This paper documents
the implemented semantic domains and their proof obligations; complete Solidity
admission and compiler preservation remain open.
\end{abstract}

\section{Research question and contributions}
We ask: \emph{How can executable Solidity semantics support mechanized
compiler-preservation proofs while retaining permitted effect orders, physical
aliasing, and an explicit distinction between source and model observations?}
The central difficulty is compositional. A call conversion may copy an aggregate after
argument side effects; a captured alias may observe a sibling write; a subsequent revert
must erase that write. Proving safe execution requires an invariant that
connects these operator rules to the heap, call frames, and failure behavior.

For example, let two initialized memory references \code{a} and \code{b}
alias the same word array, initially containing $y$ at index zero. In
\code{(a[0] = x) + b[0]}, evaluating the read before the assignment gives
$x+y$, while evaluating it afterward gives $2x$, provided the respective
addition fits in a word. Checked overflow instead produces a source panic.
A model used to justify transformations must preserve the permitted choices,
the shared physical cell and the failure behavior together.

The release provides three concrete interfaces. First, a single continuation machine
combines per-group child selection with captured physical places and an
invocation checkpoint. Second, a resolved admission boundary retains the admission
check equations; independent type, heap, and continuation invariants support
progress, preservation and resource-adequate prefix safety for a stated profile.
Third, quantified source/target pilots demonstrate calls with aliases and
rollback, and two permitted orders that can disagree on success versus overflow.
The second pilot derives machine execution and two concrete target traces for
every input word, failure flag and related entry state. An enforced theorem
manifest binds statements, axiom allowances and definition sources. Independent
compiler and EVM checks provide empirical evidence about fourteen selected interactions. Compatibility bridges connect
the certificate interfaces to the shared execution model.

Our work builds on executable Solidity semantics, memory models, and
proof-assistant verification~\cite{ksolidity,lolisa,isabelle,isabelle-shallow,memory-smt,solidcore}.
We combine independent child scheduling, physical aliasing, and rollback in
interfaces for compiler-preservation proofs. The quantified pilots test how
these interfaces compose; comparison with prior developments and external
execution checks place those results in context.

\section{Which specification is being used?}
The operational meaning of a proof in this development is fixed by the Lean
definitions in its dependency closure. The relevant specification is therefore
the executable code described here, together with the theorem's explicit
premises. Solidity documentation and Ethereum fork rules guide the source and
target definitions, whose agreement with those references remains a model-adequacy
obligation, assessed through review and differential execution.

There are three principal layers. First, \src{Contract/Language} and
\src{Contract/Semantics} define an intrinsically typed contract language.
Second, \src{Solidity/} defines the source models used by the Solidity
artifact certificates. Third, \src{Contract/Evm} gives partial bytecode
interpreters and call/creation frames. Yul word expressions and a richer
interaction model support these layers, with correspondence theorems connecting
their supported domains.

Two source boundaries must be distinguished throughout. The \emph{resolved
machine} consumes typed syntax, declaration identities and a resolved function
table; it supports a broader language, including loops and internal recursion.
The \emph{artifact certificate interfaces} relate models of a restricted
source language to particular target artifacts. Clients establish source
admission and exact artifact binding through external parsing and artifact
capture. These trusted steps connect the source files and compiler output to
the models in the theorems. The safety profile further restricts the resolved
machine to the configurations covered by its invariants. Section~\ref{sec:trust}
collects the trust assumptions for using these interfaces.

Solidity Spec is a standalone project at
\src{okcontract/solidity-spec}~\cite{artifact}. It follows our work on the
oksolc compiler. All paths in this paper are relative to this package unless
explicitly qualified. Primary declarations use the \code{Solidity} namespace;
certificate interfaces also use \code{Compiler.Validation} through
compatibility imports. The release is pinned to Solidity 0.8.36, reference
compiler identity \code{0.8.36+commit.8a079791}, Cancun, and Lean 4.29.0.
Compiler compatibility is scoped to \code{viaIR: true}, with optimization
disabled or enabled at 200 runs; legacy code generation is outside this profile.
\src{target.json} and \src{Target.lean} record those targets.
Versioned documentation provides explanatory references~\cite{control,abi,layout,types}.
Although the documentation site may label its 0.8.36 pages ``develop,'' the exact
native compiler identity pins differential tests. Accepted pragmas retain this
fixed semantic target.

\subsection{Compiler-pipeline boundary}
The choice of compiler pipeline affects the timing of implicit copies. With an initially
empty storage array \code{values}, consider:
\begin{verbatim}
values.push(1);
Holder memory h =
    Holder(values, append());
return h.a.length;
\end{verbatim}
Here \code{Holder} has a memory-array member, and \code{append()} pushes 2
onto \code{values} before returning zero for its other member. Legacy code
generation copies the first argument before evaluating the second, producing
length one. IR evaluates the raw arguments before copying, producing length
two. The corresponding internal call with a memory-array parameter behaves
the same way. Both results were reproduced by the external driver with the pinned
compiler, Cancun, and both optimizer settings. The indexed fixture returning
\code{h.a.length*256+h.a[1]} panics with code \code{0x32} under legacy and
returns 514 under IR.

The construction and invocation rules select the IR conversion phase.
Reproducing the legacy copy timing would require changing conversion placement
as well as selecting an argument order. Explicit conversion nodes could support
such an elaboration in future work. Solidity documents further differences
between the pipelines~\cite{ir-changes}, motivating the explicit compatibility
profile used here.

\begin{table*}[t]
\centering
\caption{Coverage by kind of evidence within the stated restricted
representation/admission profiles.
External checks cover selected manually paired IR cases.}
\label{tab:coverage}
\small
\begin{tabularx}{\textwidth}{@{}Xccccc@{}}
\toprule
Construct & Represented & Admitted & Generic safety & Target refinement & External check\\
\midrule
Scalar operations and word storage & Yes & Yes & Profile & Selected pilots & Selected cases\\
One-level \code{uint256[]} memory & Yes & Yes & Profile & Prepared entries & Selected cases\\
Struct construction and aggregate copying & Yes & Yes & No & No & Selected cases\\
Internal calls and reference parameters & Yes & Yes & Profile & Alias/call pilot & Selected cases\\
Events and their rollback & Yes & Yes & No & No & One rollback case\\
General source-to-resolved elaboration & No & No & No & No & No\\
\bottomrule
\end{tabularx}
\end{table*}

\section{The continuation machine and its safety contract}
\label{sec:execution}
Let $\tau$ range over resolved scalar, aggregate-reference and function types,
and let $P$ be a finite table of typed function bodies. Local cells contain typed
values or explicit uninitialized references; heap references identify physical
locations rather than copied mathematical arrays. The execution layer defines
closed resolved syntax, heterogeneous
value packets and one continuation machine. Stable local IDs and captured
physical places retain the location selected by effectful access paths. A
memory-member assignment rebinds its captured pointer cell; readback reloads
that cell. Nested tuples and holes remain groupings, without fabricating a
storable tuple value.

A configuration is $c=(t,K,s,s_0)$: a current task $t$, continuation stack
$K$, current machine state $s$, and invocation checkpoint $s_0$. The state
contains typed local cells, arithmetic mode, the physical heap and world,
external backend state, and remaining materialization work. A task may evaluate
an expression, capture a place, collect operands, invoke a function or resume
control. One transition selects a permitted child when needed and updates
these components; an invocation failure restores $s_0$.

Each eager expression group selects children independently. Independent
list-split selection rules prove that every occurrence in a completed group is evaluated once and
that every permitted permutation has a schedule; indexed results return to
syntactic order afterward. Boolean operators and conditionals are lazy.
Assignment captures its destination once. Conditional arms carry their
resolved common type; an implicit tuple-group transfer runs after the whole
selected arm has evaluated, so later argument effects precede aggregate copies.
Compound assignment permits the old read before or after the RHS. Tuple stores use reverse component order, whereas
explicit return conversions follow declaration order. These parent operations
are distinct from unspecified child-expression order.

Inline arrays allocate backing memory before evaluating their elements. Struct constructors
evaluate all raw arguments before root allocation, then convert and store
fields in declaration order. Constructor packets permit shared reference
children, allowing aliases within a constructor packet. Exiting a lexical scope
removes introduced local cells while retaining writes to outer variables.
All loop forms, break/continue and return share the same transition vocabulary.

Internal calls use a finite signature/body table and support recursion,
reference parameters, named returns and multiple results. Raw argument effects
precede parameter conversion, which precedes indirect dispatch and its
null-handle panic. Callees start checked, share allocation and the invocation
heap, and restore caller locals and arithmetic mode. Primitive store codecs
retain the certificate logical-store domains, including dirty-Boolean cleanup,
without assuming a universal physical-slot injection.

Source failures roll back the complete invocation checkpoint. Unsupported or
invalid contexts and undefined behavior are explicit inconclusive observations;
host exhaustion retains its configuration. Continuation bounds limit host
evaluation; EVM operand-stack and gas bounds are properties of a bound target
artifact. Finite soundness/completeness covers terminating runs, and an infinite
loop witness demonstrates divergence within the same semantics.

Whole-tree resolved admission checks scopes and source permissions, including
unused branches. Declaration and mutability resolution, and callee summaries
of control flow where required, are inputs to these checks. General source-byte
admission, aggregate event/error ABI encoding, arbitrary raw-memory
correspondence and compiler lowering preservation remain separate obligations.
The interfaces are defined under \src{Solidity/Core/Execution}.

\subsection{From admission to represented configurations}
Soundness and completeness of fuel-bounded iteration relate the runner to the
step semantics, whose terminal observations include inconclusive results.
Admitted-state safety additionally requires an invariant connecting syntax,
local values, captured targets, saved continuations, function handles, and
physical heap packets. The safety proof shows that transitions preserve this
invariant under explicit domain and primitive-resource assumptions.

The proved structural profile includes executable scalar types, closed recursive function
handles, canonical memory word arrays with physical aliases, scalar persistent
and transient places, eager effects, structured control and singleton
declaration prefixes. The independent \code{ConfigurationRep} predicate connects
checking equations to initialized cells, heap packets, operand partitions,
captured targets and saved continuations. These structural premises allow the
proof to derive progress. With adequate primitive resources,
\code{Machine.step} either produces a represented source observation or preserves
the invariant and checkpoint, excluding inconclusive results and primitive
suspension. Entry requires successful admission, the structural profile, a
closed admitted function table, and an initial state with represented physical
packets and closed function handles.

The profile excludes other aggregate families, storage/calldata reference
locals, tuple-pattern assignments, inline-array and struct construction,
array-length mutation, event/error encoding, assembly and linked calls.
User-defined value wrapping/unwrapping, external-function projections and
\code{delete} are also excluded; a for-loop post is an expression or skip.
\src{SAFETY.md} indexes the declarations defining this profile within the
broader executable machine.

Resource adequacy is separate from termination. Aggregate fuel is at least two
and host continuations are uncapped. Let $C=66(2^{256}-1)+96$.
The generic progress theorem assumes at least $C$ remaining host work;
an initial allowance of $(n+1)C$ suffices for an $n$-transition prefix.
These conservative bounds suffice for the generic proof; particular executions
can use less work. The extra allowance covers
the runner's real terminal inspection even at zero transition fuel. The
\code{execute\_prefix} theorem covers every child-choice sequence. A completed
prefix has represented values or exact source-failure rollback; an exhausted
prefix retains the invariant, original checkpoint and a further $C$ allowance.
The bound measures modeled materialization work; EVM gas and physical host
capacity require separate analyses. Final packet typing is existential in the
terminal predicate, so the generic theorem establishes representation
preservation and absence of model failure. Relating those packet types to the
root expression's result type remains open. The pilot theorems establish the
exact returned words for their particular programs.

Two regression families illustrate the distinction. Invocation performs
the permitted implicit internal-function conversion, including pure-to-view
widening, while retaining the handle. Separately, eager siblings receive the
same incoming reference-readiness facts. Assignments guaranteed by a complete
group become available afterward; a read requiring another sibling to run
first is unsupported by this defined admission domain. Lazy operands and
sequential statements retain their required order. The raw machine still
represents the original order-dependent inconclusive schedules.

The lightweight import \src{Solidity.Execution} exposes checked entries,
the machine and its execution facade. Additional execution laws are available
through \src{Solidity.Execution.Proofs}; admitted-state safety is imported
through \src{Solidity.Execution.Safety}. Compatibility clients use
\src{Solidity.Compatibility.Execution}; the comprehensive \src{Solidity}
umbrella includes all four interfaces. Indexed checked entries retain
the admission equations. Their projection proofs preserve the facade,
including failed preparation. \src{Policy.lean} defines fixed compiler choices,
with source provenance recorded in \src{references/policies.json}. These
interfaces share the same normalizer and interpreter.

\section{Observation contract and provenance}
We distinguish three observation profiles. A \emph{functional} observation
compares completed return/revert data and the selected public state and effects.
A \emph{metered} observation additionally compares the specified resource
outcome. A \emph{representation-sensitive} observation also binds the exact
code or packet representation required by its relation. These profiles are
recorded in \src{SEMANTIC_CONTRACT.md}, with explicit premises for each relation.
\src{Observation.lean} defines completed records and named
relations over invocation outcomes, public state, transient storage,
account lifecycle and ordered interactions. Stronger profiles also compare
resource observations and the bound code representation. Correspondence requires
completed, supported execution and proof of any capability restrictions stated
in the summary.

Invocation rollback restores the entry snapshot, including logs and
transaction-local storage. Gas consumption and transaction admission/fee
settlement have separate boundaries. The current profile starts after
admission and intrinsic gas accounting, with fee settlement handled externally.
Account existence and Cancun creation/deletion facts are tracked explicitly
alongside balances and code.

Sixty-two reference snapshots are pinned by immutable upstream revision and
content digest in \src{references/lock.json}. They include selected Solidity
documentation and type rules and the Cancun execution-specs fork directory.
The lock distinguishes complete snapshots from documentation indexes with
unbundled chapters, making the reference corpus inspectable at a fixed revision.

\section{Quantified source/target experiments}
\subsection{Internal calls, aliases and rollback}
The composed pilot quantifies over input words $x$, failure flags
$b$ and related initial states. Two initialized memory references alias one
represented length-one word array. An ordinary entry call invokes an
internal helper, which writes the array payload and persistent slot zero. The
caller reads through the other alias and either returns $x$ or reverts. The
source theorem derives parameter binding, named-result initialization, both
calls, result transfer and caller restoration through the continuation machine.
It establishes a checked completion within 160 transitions with choices
\code{[]} for every materialization allowance.

The target is a manually paired, hand-authored 23-byte Cancun fragment. Its
prepared-entry relation gives six active memory words and zero bytes beyond
that extent. A metered target trace determines adequate-gas observations;
failure and insufficient gas restore the complete target world and transaction
snapshot. Frame theorems preserve other owners, other storage slots and memory
packets disjoint from the written payload. The proof derives both the source
run and its correspondence with this trace from the entry-representation premises.

\code{refines\_represented} composes this same budgeted source execution with
generic safety at $B\geq161C$, then derives matching adequate-gas target
observations and separate insufficient-gas rollback. Both exported refinements
retain $g_{\rm remaining}=g-\mathit{used}$ for every
$g\geq\mathit{needed}$, with $\mathit{used}\leq\mathit{needed}$.
Their source/target relation is functional; the equation additionally specifies
target gas consumption, without assigning EVM gas to the source. A nonemptiness theorem
constructs its represented, related domain for arbitrary persistent worlds,
owners, words and flags. A negative fixture shows that the default 4096 work
allowance suffices for this nonallocating trace but falls below the generic
uniform bound. A companion applies all-choice prefix safety to the
effectful aliased group \code{(a[0] = x) + b[0]}, permitting order-dependent
values and source overflow panic.

Both experiments start from prepared entry states and hand-authored artifacts.
Their scope is execution from those entries; source-byte correspondence,
compiler emission, and deployment/allocation prefixes remain external to the
composition. \src{PILOT.md} indexes the proof modules and reproduction commands.

\subsection{Order-sensitive success and rollback}
The second experiment uses two initialized references to the same initially
zero memory word and the following prepared function body:
\begin{verbatim}
stored = x;
uint256 z = (a[0] = x) + b[0];
if (fail) revert();
return z;
\end{verbatim}
The source has two concrete permitted schedules. Reading the alias first gives
$x$; writing first gives $2x$ if it fits, and panic \code{0x11} otherwise.
For $x=2^{255}$ and a false failure flag, one schedule returns $2^{255}$ and the
other restores the complete entry state. The initial storage write therefore
participates in the rollback argument.

For each schedule, the \code{source\_execution} theorem derives machine execution
with a transition allowance of 80 for every input word, failure flag and host-work
allowance, including zero. The \code{source\_permitted} theorem embeds that
trace in the permitted-execution relation. Admission is checked for the
prepared body and its expected word return binder. Independent premises require
represented aliases, an initially zero payload and writable storage. A domain
witness establishes these premises and related target states for arbitrary
persistent worlds.

Two fixed, hand-authored Cancun fragments implement these schedules: 25 bytes
for the read-first case and 111 bytes for the write-first case. Both receive the
original word and flag and perform the storage write. The latter checks the
high-half bound before doubling. Their guarded traces determine a gas
threshold and consumption, including the storage sentry. The exported
\code{refines} theorem proves functional correspondence for every adequate gas
budget, the exact remaining-gas equation, and complete target world/transaction
rollback for overflow, explicit revert and insufficient gas.

Negative controls reject a wrapping zero return at $2^{255}$ for both contracted
schedules and reject a target that matches the empty revert bytes while retaining its
write. These results compose order, aliasing and failure in one source/target
argument. The body operates on existing aliases; internal calls are covered by
the first pilot and aggregate deep copies by the external tests. Its
allocation-free trace also gives a zero-work bound for this program, illustrating
how a program-specific argument can sharpen the generic profile's uniform bound.

\section{Released proofs and executable checks}
\label{sec:evaluation}
The release separates proof checking, model examples and external conformance.
The package build elaborates both libraries and their regression modules. Build checks use four Lean workers. Fresh kernel
replay checks the comprehensive library, the call/resource, child-order and order-sensitive pilot roots and the scalar
machine regression closure, including their imports. These checks concern the released Lean definitions and proofs. External
compiler execution is reported separately below.

\subsection{Theorem contracts and composition}
\src{theorems.json} records 85 selected roots with fully qualified names,
expected elaborated statements, complete kernel-statement hashes and permitted
axioms. A statement hash binds references to named predicates. The manifest additionally
records SHA-256 identities for the imported project definition closure, build
inputs and auditor. The release command compares compiled statements,
complete kernel-statement hashes, per-root transitive axiom allowances and this
source closure before fresh replay. Changed definitions therefore fail the gate
even if a theorem's displayed statement is unchanged. Recording new expectations
is an explicit review operation and is never an acceptance run. Fresh replay
then checks proof terms independently of elaboration.

The two pilots compose checked admission and independent representation
premises with source execution, target execution, gas conditions, and rollback
for all quantified inputs. Explicit bounds on source transitions, aggregate
recursion, and materialization work ensure that these executions complete.

\subsection{Model fixtures}
The package includes manually paired and generated model fixtures under
\src{Solidity/Conformance}. The scalar corpus contains 45 scalar, loop and
constant cases and ten constant rejection examples. The inference corpus
contains 38 expression cases and 16 rejection examples; the static matrix
contains 17,690 generated conversion, operator and update cases. These modules
provide executable model expectations and source representations.

The aggregate corpus contains 46 runtime examples, five static rejection
examples and nine model-boundary checks. The effectful corpus contains 48
machine programs, three source-rejection examples and two host-boundary checks.
The interaction corpus contains 42 generated combinations of aliasing,
evaluation order, conversions, calls and rollback. Expected outcomes are
computed by executing the model, with manually paired source representations.
These full-corpus counts describe model fixtures; the following external
evaluation independently executes a selected subset.

Kernel regression modules exercise ABI canonicality, payability, storage
layout, stateful execution, constructors, admission rejection, rollback and
host exhaustion. Import and declaration-dependency guards check that primary
execution interfaces use the shared machine. These concrete regressions
complement the universal theorem contracts and help localize changes in behavior.

\subsection{Independent compiler/EVM evaluation}
The standalone driver executes fourteen selected effect cases through the
Lean machine and checks resolved admission, then compiles the paired Solidity
with the exact reference compiler and runs the artifacts on a local Cancun
Anvil instance. Both IR optimizer settings are checked. The 28 observations
matched exported permitted outcomes: 20 returned and eight reverted. There
were no unexpected compiler rejections, model admission failures, unsupported
executions, host exhaustions or observation discrepancies in this runtime set.

The comparison tests membership in a finite set of executed, permitted schedule
witnesses. Because this set may cover only some schedules, a nonmember requires
examining the pairing and coverage before diagnosing a semantic mismatch.
The cases include aliased
operand ordering, argument effects and copies, compound destination capture,
overlapping tuple stores, internal reference calls, checked overflow, and
rollback. Each call observes unchanged initial contract storage. Reverting
cases additionally execute as transactions with failed receipts, empty logs
and unchanged selected storage slots. The rollback check observes precisely
these components of the transaction result.

Three separate static-rejection controls matched compiler diagnostic identities.
Two model resource controls preserve exhaustion as a host limitation. The three
copy-timing probes also run under legacy compilation with both optimizer
settings: all six observations exhibited the documented boundary discrepancy
and are recorded separately as negative compatibility controls.
The runner retains exact source and Standard JSON inputs, compiler outputs,
tool identities/hashes, RPC observations, process statuses and package input
identities, allowing the corpus observations and compiler configuration to be
reproduced. The manual source/model pairing remains a review obligation.

\section{Domains and semantic layers}
Let $M=2^{256}$, $\W=\{0,\ldots,M-1\}$, $\mathbb A=\{0,\ldots,2^{160}-1\}$,
and $\mathbb B=\{0,\ldots,255\}$. Lean represents words and addresses using
\code{BitVec 256} and \code{BitVec 160}. Bytes are \code{UInt8}; byte strings
are finite lists. We write $\rho$ for a local environment, $\sigma$ for storage,
and $\kappa$ for context.



\subsection{Word operations}
\src{Word.lean} defines wrapping word operations and explicit
checked operations. Wrapping addition, subtraction and multiplication operate
modulo $M$. Logical shifts give zero for a shift count at least 256. The shared
low-level unsigned division and remainder follow the EVM convention of yielding
zero at a zero divisor; the checked variants below produce panic bytes.

Define $P(c)$ as the four selector bytes \code{4e 48 7b 71}, followed by 31
zero bytes and the one-byte code $c$. The checked arithmetic evaluator returns
\code{Except (List UInt8) (BitVec 256)}. For natural representatives $a,b$,
\begin{align*}
 \operatorname{add}_{c}(a,b)&=
 \begin{cases}\operatorname{ok}(a+b)&a+b<M,\\P(\code{11})&\text{otherwise},\end{cases}\\
 \operatorname{sub}_{c}(a,b)&=
 \begin{cases}\operatorname{ok}(a-b)&b\leq a,\\P(\code{11})&\text{otherwise},\end{cases}\\
 \operatorname{mul}_{c}(a,b)&=
 \begin{cases}\operatorname{ok}(ab)&ab<M,\\P(\code{11})&\text{otherwise}.\end{cases}
\end{align*}
Here panic codes are hexadecimal and the second branch denotes an error.
Checked division and remainder use $P(\code{12})$ at
zero. Assertion failure in the typed language uses $P(\code{01})$. These choices
follow the intended checked-arithmetic and panic behavior~\cite{control}.
Only addition, subtraction and multiplication from this checked family occur
in the admitted Solidity word-expression grammar.

\subsection{The typed contract language}
\src{Basic.lean} defines types
$\tau::=\mathrm{word}\mid\mathrm{bool}\mid\mathrm{address}$.
A variable is indexed by a type list $\Gamma$ and its type; the matching
\code{Env} stores values of exactly those types. Storage with $n$ declared
slots is a function $\mathrm{Fin}(n)\to\W$. Consequently, unbound variables,
out-of-layout accesses, and ill-typed operands are unrepresentable in this AST.

Expressions include literals, variables, storage reads, caller/self addresses,
word operators, comparisons and negation. Statements include returns, explicit
reverts, local binding and assignment, checked operations, assertions, stores,
branches, finite helper calls and event emission. Helpers are embedded closed
bodies, giving a finite, nonrecursive call structure. Event topic count is
bounded by four with a proof field.

The evaluator in \src{Eval.lean} is a total recursive
function with the big-step interpretation
\[
 \eval{s}(\sigma,\kappa,\rho)
 \in \operatorname{Success}(\sigma',\vec v,L)
       +\operatorname{Revert}(b).
\]
Stores update one slot before evaluating the continuation. Helper results are
prepended to the caller's environment; helper storage changes feed the
continuation. Logs are ordered by execution. Revert carries no post-state and
no retained logs. At the invocation boundary,
\[
 \operatorname{commit}_{\sigma}(o)=
 \begin{cases}\sigma'&o=\operatorname{Success}(\sigma',\vec v,L),\\
 \sigma&o=\operatorname{Revert}(b).
 \end{cases}
\]
Thus a store followed by revert has the same committed storage as immediate
revert. The theorems \code{revert\_rolls\_back} and
\code{store\_then\_revert} express this directly. Gas, ABI entry, and deployment
are outside this typed evaluator. In particular, plain expression addition is
wrapping; checked behavior is an explicit statement constructor.

\section{The shared typed core}
\subsection{Values, operators and permitted orders}
\src{Scalar.lean} indexes integer and fixed-byte values by
$w\in\{8,16,\ldots,256\}$. Unsigned and signed integers share a bit-vector
representation but have distinct types and operations. Signed interpretation
uses two's complement. Checked addition, subtraction and multiplication test
the corresponding mathematical range; unchecked operations wrap modulo $2^w$.
Signed division truncates toward zero. The minimum signed value divided by
$-1$ panics in checked mode and wraps in unchecked mode. Division and remainder
by zero always emit $\operatorname{Panic}(\mathtt{0x12})$, including within
unchecked expressions. Arithmetic overflow emits $\mathtt{0x11}$.
Shifts truncate without overflow checks; signed right shift sign-extends.
The represented casts truncate or extend according to their typed constructor;
these primitive operations are separate from source conversion permissions
and literal elaboration.

Arithmetic mode belongs to each operator, so resolved syntax can represent
nested checked/unchecked regions. \src{Expression.lean}, relative to
\src{Solidity/}, defines literals, typed references, unary/binary operators,
conditionals, conjunction and disjunction. References come from a resolved
program. The executable evaluator chooses a uniform operand order.
The separate relation \code{Expr.Evaluates} permits each binary node to choose
independently; \code{eval\_permitted} proves executable evaluation sound for that
relation. The relation includes outcomes unavailable under either uniform
strategy alone. Lazy constructs evaluate only the necessary operand/branch.
These rules describe the retained pure fragment. The effectful machine in
Section~\ref{sec:execution} connects its operand strategies to a shared heap
and continuation stack.

\subsection{Correspondence and stateful consumers}
The modules \src{Word.lean},
\src{TypedExpressions.lean}, \src{Conditions.lean}, and \src{Storage.lean}
(the latter three in the same directory) prove embeddings for arbitrary
expressions, environments and current stateful source statement bodies.
The typed DSL's zero-totalized EVM division is translated to a conditional that
preserves its zero-divisor result within the shared core.

\src{Statement.lean} defines finite blocks over a typed store
interface. The concrete compatibility store reuses source-derived field layout,
ABI locals and static-write checks. Sequential composition propagates a terminal
outcome; reverted and exceptional outcomes expose no committed store.
\code{body\_result} equates this evaluator with the certificate stateful evaluator.
\code{certified\_contract} then proves that a complete call certificate
implies \code{TracePreservesCore}: the shared-core result conforms to the exact
target trace for every sized input, store and static-mode setting, retaining
the certificate's source domain and trace-derived gas observations.
Typed helper/log statements and specialized ABI/constructor clients have
explicit shared-core correspondence, extended below to the effectful machine.

\subsection{Typed frames and certificate clients}
The procedure layer reuses core expressions, environments and stores for
finite closed helpers, local assignment, tuple returns and static-word logs.
The theorem \code{Compatibility.Typed.procedure\_eval} covers every statement
of the retained typed language. Helper results extend the caller's frame;
a later failure discards both helper and caller logs and writes.
The resolved effectful machine adds recursion; linked external execution remains
a separate obligation.

Ordered declaration initializers and complete constructors execute core
blocks, with proved outcome correspondence. Constructor bare returns retain
their initialization meaning. Specialized ABI admission is related to the
shared canonical-word predicate; empty, raw-word, literal and expression
methods have core outcome adapters. Within the existing source domain, complete
certificate transfers retain the calldata, static-mode, rollback, deployment
and gas claims.

\subsection{Structured control and lexical resolution}
\src{Control.lean} adds a continuation machine for sequencing,
branches, loops, break and continue. A loop carries a finite post block: skip
for a while loop, or the increment for a for loop. Continue runs that post block,
while break skips it. A return or revert bypasses all enclosing continuations.
Finite interpreter execution is connected to a fuel-free transition judgment:
\[
 \operatorname{run}(k,c)=\operatorname{finished}(o)
 \quad\Longrightarrow\quad c\Downarrow o,
\]
and every finite derivation has a sufficient interpreter budget. Separate
theorems prove preservation under additional fuel and determinism for a fixed
operand strategy. Exhaustion records the reached configuration as a host outcome,
allowing finite-prefix reasoning for programs whose loops may diverge. Connecting
these resolved loops to source text and target gas requires elaboration and
artifact proofs.

\src{Resolve.lean} checks lexical references in typed scalar
expressions and constructs de Bruijn variables. The nearest declaration wins;
a type mismatch at that declaration is rejected rather than searching past it.
Unknown names are rejected even in a branch skipped at runtime. Soundness and
completeness relate the checker to an independent structural elaboration
judgment. Evaluation preservation assumes agreement on resolved bindings.
Parsing, type inference, overloads, inheritance and exact source-byte
correspondence remain future frontend obligations.

\subsection{Exact constants and checked preparation}
\src{Literal.lean} evaluates constant-expression trees over
normalized rational numbers. Addition, subtraction, multiplication, division,
remainder, exponentiation, bit operations and shifts follow the pinned 0.8.36
rules. Remainder uses a quotient truncated toward zero; right shift of a
negative integer rounds downward. Division by zero, fractional bit operands,
invalid shifts, precision exhaustion and non-representable conversions produce
static elaboration errors. A universal bit-observation theorem
proves the Boolean meaning of signed constant bit operations, including every
sign-extension bit.

The result precision guard checks the highest set-bit index of the normalized
numerator and denominator against 4096. Exponentiation and left shift also
apply conservative work guards. Thus $2^{2049}$ is rejected as a literal power
even though its result would satisfy the separate result guard. The pinned
compiler treats $0^e$ as zero for every nonzero literal exponent, including
negative exponents. Such edge cases are explicit versioned rules.

Integer conversion is characterized in both directions. For unsigned width
$w$, with $q$ a normalized rational and $v$ a $w$-bit word,
\[
 \begin{aligned}
 &\operatorname{toUnsigned}_w(q)=\operatorname{ok}(v)\\
 &\quad\iff q.\mathrm{den}=1\;\land\;\operatorname{toNat}(v)=q.\mathrm{num}.
 \end{aligned}
\]
The signed theorem replaces the final observation by
$\operatorname{toInt}(v)$. Both characterizations require exact representability
of the constant, whereas runtime narrowing casts truncate to the target width.

\src{Traversal.lean} implements checked substitution of expression
leaves and its structural soundness/completeness theorems.
\src{Prepare.lean} uses it to resolve names and convert exact
constant subexpressions. Its separate elaboration judgment requires lexical
binding or integral, exactly representable constant values. The checker is
sound and complete for that judgment and preserves evaluation when source
leaves have their specified meanings. It checks all branches, including those
skipped at runtime, operating on parsed trees with explicit operator and type
annotations.

Normalization precedes the result-precision check. For example, adding two
copies of $1/2^{2049}$ has a raw denominator $2^{4098}$ but a reduced denominator
$2^{2048}$. The reduced rational satisfies the result-precision guard, as
demonstrated by kernel-checked examples in the released literal fixtures.

\section{The admitted Solidity source semantics}
\subsection{Resolved pure word expressions}
The syntax in \src{Word.lean} is
\[
 e ::= n\mid x\mid {\sim}e\mid e\mathbin{\circ}e,
 \qquad
 \circ\in\{+,-,*,\&,|,\mathbin{\oplus},\ll,\gg\}.
\]
\code{Expr.eval} maps a name environment $\rho:\mathrm{String}\to\W$ to a
word, with literal conversion modulo $M$.
\src{Checked.lean} lifts evaluation into an exception monad,
checking arithmetic and propagating exact panic bytes. Bitwise and shift
operators retain their wrapping/truncating meaning~\cite{types}.
The source model uses one checked flag for the entire body.

The AST is deliberately more permissive than the admitted source text. A
natural literal in Lean can exceed the word range, and the environment is
total. The trusted parser must establish literal bounds, name resolution,
types and allowed constructs. It rejects literal arithmetic with Solidity's
unbounded rational-literal semantics; operands instead require admitted word
values such as \code{uint256(3)}. This prevents silently interpreting a
rational constant expression as modular arithmetic.

Checked expression evaluation visits operands left to right. In this pure
fragment, all arithmetic failures use the same panic bytes, making the eager
evaluation order unobservable. Boolean short-circuiting in the stateful model
is explicit and does affect whether an operand can panic~\cite{control}.

\subsection{Function inputs and observations}
\src{Function.lean} records a name, word parameter names, body,
checked flag and payability flag. Input consists of calldata $d$ and incoming
value $v$. Parameter $i$ is decoded from byte offset $4+32i$ as a big-endian
word. Successful word results are encoded as 32 bytes; arithmetic failure
returns the panic bytes as a reverted outcome.

A signature's selector is the first four bytes of Keccak-256 over its canonical
name and parameter types~\cite{abi}; the implementation is
\src{Keccak.lean}. The stronger call interface in
\src{Call.lean} accepts a word method exactly when
\begin{align*}
 A_f(d,v)\ \Longleftrightarrow\;&(f.\mathit{payable}\lor v=0)\\
 &\land |d|\geq 4+32\,\mathit{arity}(f)\\
 &\land \operatorname{selector}(d)=\operatorname{selector}(f).
\end{align*}
Failure of admission gives empty revert bytes. Trailing calldata is allowed.
The separate \code{ValidCall} predicate requires exact canonical length and
zero value. It witnesses inhabitation within the broader input domain of the
all-calls certificate.

\subsection{Typed ABI and dispatch}
\src{ABI.lean} adds word, Boolean and \code{bytes4} parameters.
A Boolean word must be zero or one; fixed bytes occupy the high-order bytes,
with zero padding in the remaining 28 bytes~\cite{abi}. In notation,
\begin{align*}
 A_{\mathrm{bool}}(w)&\iff w\in\{0,1\},\\
 A_{\mathrm{bytes4}}(w)&\iff w\mathbin{\&}(2^{256}-2^{224})=w.
\end{align*}
Malformed complete argument words therefore revert during decoding.

The contract-level interfaces compose complete runtime-entry traces for
short calldata, unknown selectors, disallowed value, truncated arguments,
malformed arguments and selected body results. Selector uniqueness is a
certificate condition. Separate models cover empty bodies, exact raw-word
returns, literal Boolean/address/\code{bytes32} returns and empty contracts,
with expression support limited to each model's grammar. The stateful source
model also supports an optional receive body for empty calldata; general
fallback handlers and external calls remain beyond this fragment.

\subsection{Stateful statements and control}
\src{StorageSource.lean} defines comparisons, Boolean conditions,
statements and functions over a loop-free grammar. Statements assign a word or
Boolean field, return a word, stop with empty return data, revert, or branch
into finite statement lists. Conditions have literals, Boolean field reads,
unsigned comparisons, negation, conjunction and disjunction.

Evaluation has an explicit continuation result:
\[
 \operatorname{Flow}=\operatorname{Next}(\sigma)
 +\operatorname{Halt}(\operatorname{Outcome}),
\]
where an outcome is returned bytes plus storage, reverted bytes, or a fault.
The first terminal result skips the rest of the statement list. Conjunction
and disjunction short-circuit; a panic in a required condition propagates before
either branch executes. Assignment evaluates its right-hand side first, then
checks static mode, then writes the declared field. Invalid declaration indices
fault in the evaluator and are rejected by source admission. Falling through
a word-returning body returns the zero word; a void body returns empty bytes.

\section{Storage layout and initialization}
\src{StorageLayout.lean} derives layout directly from declaration types.
The cursor is a slot number and a byte position in $\{0,\ldots,31\}$.
Boolean fields consume one byte from
the low-order end. A word field occupies an entire slot; if the cursor is
partway through a slot it advances to a fresh one~\cite{layout}.
For example, declarations \code{bool a; bool b; uint256 c;} occupy
$(0,0),(0,1),(1,0)$ respectively.

Let $m_j=255\ll 8j$. Boolean extraction tests whether the selected byte is
nonzero and returns a canonical zero/one word. Writing Boolean $b$ uses
\[
 \operatorname{insert}(w,j,b)=
 (w\mathbin{\&}{\sim}m_j)\mathbin{|}
 (\operatorname{encode}(b)\ll8j).
\]
The theorem \src{PackedBool.insert_preserves_other_bits} preserves every
bit outside that byte, including arbitrary unused bits. Whole storage is a
function $\W\to\W$, so certificate statements can quantify over arbitrary
initial stores. The source parser's layout domain is private word/Boolean
declarations; mappings, arrays, inheritance, custom layout, and synthesized
public getters require broader source models.

\src{StorageInitialization.lean} evaluates declaration initializers
in order, starting with zero storage. Each word initializer uses checked
arithmetic and sees preceding writes; a forward read sees the field's current
(default-zero) value. Boolean declaration initializers are currently rejected
by the source parser even though the model has Boolean storage operations.

\src{StorageConstructor.lean} composes these initializers with an
optional parameterless constructor body. Nonpayable constructors reject
incoming value. The body has its own checked flag, but initializers remain
checked even if the body is unchecked. A bare return completes initialization
with the current storage, while a source revert or fault preserves its failure
classification.

\section{Target semantics and gas}
\subsection{Partial raw-byte execution}
\src{Subset.lean} defines a bytecode machine with program counter,
word stack, byte memory and allocated-memory word count. Instruction decoding
reads raw bytes; jump destinations exclude PUSH data. Stack size is bounded
by 1024. Missing code decodes as STOP. The admitted read-only opcodes include
word arithmetic, comparisons, bitwise operations, shifts, stack operations,
calldata/value access, code copy, memory operations, jumps and terminal
instructions. Other opcodes return an explicit unsupported result.

The memory expansion charge for $w$ words is
\[
 C_{\rm mem}(w)=3w+\lfloor w^2/512\rfloor.
\]
Execution has separate outcomes for a completed trace, fault, unsupported
opcode and exhausted fuel. Interpreter fuel bounds the trace length inspected
by the proof, while gas accounting applies the target's execution charges.
\src{Execution.lean} proves that a completed unmetered execution
with cost $c$ determines execution at every gas budget:
\[
 \operatorname{meter}(c,o,g)=
 \begin{cases}\operatorname{finished}(g-c,o)&c\leq g,\\
 \operatorname{exceptional}&g<c.
 \end{cases}
\]
Larger step bounds preserve the completed result. Gas costs belong to each
specific target artifact and can differ across optimized implementations.

\subsection{Storage execution}
\src{Storage.lean} extends the same decoder with SLOAD and SSTORE.
Its state contains current storage, warm keys and a signed refund counter;
its environment records original storage and static mode. With original,
current and new values $o,c,n$, the SSTORE charge is
\begin{align*}
 &(\text{warm? }0:2100)\\
 &\quad+\begin{cases}
 100&c=n,\\
 20000&o=c=0,\ c\ne n,\\
 2900&o=c\ne0,\ c\ne n,\\
 100&o\ne c,\ c\ne n.
 \end{cases}
\end{align*}
The model applies the Cancun-era warm/cold schedule and refund adjustments;
transaction refund settlement and its cap are outside the execution frame.
The relevant rule references are EIP-2200, EIP-2929 and EIP-3529~\cite{eip2200,eip2929,eip3529}.

SSTORE requires remaining gas strictly above 2300 even when consumption is
less. Consequently, \src{Guarded.lean} records both gas used $u$
and minimum initial gas needed $h$. Charging a step of cost $c$ and guard $q$
to a suffix $(u,h)$ yields
\[
 (c+u,\;\max(q,c+h)).
\]
The metering theorem tests $h\leq g$ and returns $g-u$; testing only total
consumption would be insufficient. Static SSTORE faults. Frame observation
commits successful storage, warming and refunds, while revert/exception retain
the pre-frame persistent world and discard reverted effects.

\subsection{Call and creation frames}
\src{Frames.lean} and \src{StorageFrames.lean}
connect machine results to account worlds. A successful call commits incoming
value and permitted state changes. A revert retains its bytes and unused gas
but rolls back value/state; an exception has empty bytes and no remaining gas.
Frame certificates require completed, supported execution.

Creation certificates bind exact creation and runtime bytes, execute
initialization and install the certified runtime. They account for execution
and code-deposit gas, runtime size and prefix checks, and the initcode size
bound~\cite{eip170,eip3541,eip3860}. They then derive runtime-call preservation
at the deployed account. Their gas budget begins after transaction admission
and intrinsic/initcode charges. Address selection, creator nonce admission,
collision checks, and transaction fees belong to the external transaction
boundary. The interaction model separately represents additional creation
operations.

\section{What the certificates prove}
Let $S$ be an admitted source model, $b$ exact runtime bytes, $k$ a proved
step bound and $i=(d,v)$ an input. The pure-function interface has the shape
\begin{align*}
 &\operatorname{Preserves}(S,b,k)\ \Longrightarrow\\
 &\quad\forall i,\ |d|<2^{64}\ \Longrightarrow\ \exists c,\\
 &\qquad\operatorname{run}(b,i,k)=
       \operatorname{finished}(c,\operatorname{outcome}_S(i)),
\end{align*}
together with an inhabited canonical-call domain, the metering result for
every gas budget and larger fuel, and preservation of the account world under
an installed-code premise. The calldata size bound restricts the certificate's
proof domain independently of EVM consensus rules. Domain inhabitation concerns
admitted inputs and therefore also covers functions that always panic.

The storage endpoint additionally quantifies over original/current storage,
warmth, refunds, static mode and every gas budget. Its source outcome fixes
return/revert bytes and complete committed storage; the exact target trace
determines execution costs, access effects and refunds. Contract endpoints
compose all dispatch paths from bytecode entry. Deployment composition discharges the
installed-code premise by deriving the installed image from exact creation
execution. Stateful creation requires a successful deployment witness
conditional on successful source initialization, allowing always-reverting
constructors within the domain.

These implications apply to source/artifact pairs accepted with a complete
proof. A sound validator may reject correct compiler output or stop because of
unsupported instructions or resource limits. Compiler-specific lowering and
artifact capture are external clients responsible for producing the pairs and
discharging the certificate premises.

\section{Trust and exclusions}
\label{sec:trust}
Certificate validity depends on four explicit boundaries:
\begin{enumerate}
\item \emph{Source admission.} Certificate clients must relate complete source
text to the formal models. Token recognition, name/type resolution, pragma
admission and source-to-model correspondence currently rely on an external,
unverified frontend.
\item \emph{Artifact binding.} A certificate client binds exact source names/content,
settings, creation bytes and runtime bytes. Completeness and fidelity of the
native artifact capture remain obligations of that client.
\item \emph{Model adequacy.} The Lean evaluators must faithfully describe the
intended Solidity fragment and Cancun behavior. Agreement with external
specifications is supported by source review and differential evidence;
a complete mechanized correspondence remains open.
\item \emph{Proof checking.} The pinned Lean kernel checks proof terms.
The selected theorem contracts record permitted uses of \code{propext},
\code{Classical.choice} and \code{Quot.sound}. The release combines fresh
imported-kernel replay with a separate manifest comparison. Issuing an artifact
certificate additionally requires a proof for the particular source/artifact pair.
\end{enumerate}

The restricted source language of the certificate models excludes general signed/narrow arithmetic,
division and remainder expressions, loops, recursion, external calls,
inheritance, modifiers, mappings, dynamic ABI data and arbitrary assembly;
assembly support is confined to specific raw-return forms. These restrictions
apply to the certificate endpoints even where the resolved machine or
interaction model supports broader operations. External admission tools may
impose further resource limits as implementation policies.

The interaction evaluator uses an explicit call oracle and a creation executor.
Its environmental reasoning is conditional on their contracts; verifying the
called bytecode and connecting it to those contracts remains a client obligation.


\section{Resolved scalar semantics and inference}
The resolved scalar interface retains source distinctions erased by machine words:
payable addresses, enum identity and cardinality, value-type identity and
underlying elementary type, contract ancestry, and function visibility,
mutability and argument/result types. A resolved declaration registry has an
independent well-formedness judgment and an equivalent executable checker.
Registry construction and declaration resolution remain later obligations.
Aggregate, tuple, metatype and fixed-point forms are represented explicitly.
The aggregate operations below consume the resolved registry; tuple effects,
metatype members and fixed-point execution retain their separate boundaries.

Independent conversion judgments separate implicit conversion, explicit
source cast syntax and internal machine coercion. Integers widen exactly and
narrow modulo the target width; fixed bytes resize on the opposite end.
Literal conversion requires representability. An out-of-range runtime enum
conversion panics with code 0x21, whereas an out-of-range literal conversion
rejects statically. Function arguments and results are invariant under implicit
conversion, with the specified mutability weakening. Internal function handles
may identify compiler-equivalent declarations, so the model permits shared
code identities.

The unannotated expression tree records names, literal spellings and source
operators. Inference selects operand types and inserts permitted conversions.
For example, an expression \code{x + 256} with \code{uint8 x} obtains
\code{uint16}; \code{x + 1} retains \code{uint8}. Both conditional arms and
logical operands are checked before lazy execution. Thus an invalid cast in
an unreachable arm still rejects. Checker soundness/completeness is relative
to the independent resolved typing judgments over these expression trees.

Binary exponentiation shares the core's width-indexed words. Its wrapping
algorithm agrees with mathematical power modulo the word modulus; a saturated
integer algorithm implements the checked bounds without constructing a huge
power. Universal signed/unsigned laws and correspondence to the independent
target power model are proved. Large shift counts use guarded scalar shifts.
The pure execution relation permits independent ordinary operand orders; the
strategy realizes that relation. A progress theorem excludes static-rejection
and unsupported markers from typed expression execution: the result is a value
or a Solidity panic. A target-resource claim additionally requires artifact
refinement.

Literal semantics covers denominations, scientific notation, numeric hex,
address checksums and ASCII/Unicode/hex strings. String UTF-8 checks follow
adjacent-token concatenation. Pinned upstream quirks are explicit: a nonzero
\code{1e+2} rejects, \code{0e+2} is zero, and a leading-zero exponent is parsed
as octal, making \code{1e010} equal to $10^8$. Mantissas are normalized before
the scientific precision precheck. \src{SCALARS.md} describes these
version-specific compatibility rules. The exact-integer
reification of the binary64 precheck and concrete Keccak implementation retain
their stated model-adequacy obligations.

Scalar assignment/update plans and pure free-function operator signatures
guide location, effect and declaration checking. The continuation machine
executes the corresponding lvalue effects and function bodies, while global
overload resolution, full parsing, and source-byte binding remain frontend
obligations. \src{SCALARS.md} indexes the scalar interfaces and theorem families.

\section{Aggregate locations and representations}
The aggregate interfaces reuse the shared interaction world, byte memory,
width-indexed values and concrete Keccak implementation. Persistent and
transient account stores are distinct. Arrays, structs, mappings, bytes,
strings and calldata slices carry physical references with invocation ownership,
packed byte positions, slice descriptors and captured memory binding cells.
Reference values clear the former binding cell when read as rvalues.

Layout computes packing, static extents, inherited declaration order and custom
persistent bases, with a separate transient cursor. The checked declaration
endpoint validates every physical placement, preserves the complete retained
declaration inventory and proves bounded cursors and fitting packed fields.
The linearization and resolved registry are inputs from declaration elaboration.
Runtime storage strides and addresses wrap as 256-bit words. The representation
relations allow distinct logical paths to share physical locations, including
collisions in Keccak-derived addresses.

Local storage and calldata references alias; memory-to-memory reference
assignment also aliases and rebinds a captured member cell. Storage-target
copies execute ordered physical operations, including resize-before-copy,
buffered packed-word reads and byte cleanup/header order. Independent recursive
judgments characterize successful execution, and exact final-store folds retain
collisions. Produced scalar observations compose across later disjoint bit
fields as well as separate slots. For overlapping copies, the final-store folds
account for reads that can observe earlier writes in the same copy. Fresh-memory
copies collect immutable-source packets with source-specific decoding/allocation
check order and populate new
objects. Constructor induction proves scalar words, pointer cells, headers and
payloads, while staged constructor theorems derive freshness from the
intermediate allocation frontier.

Defaults permit the canonical shared zero-slot object for empty dynamic values;
fixed arrays and structs are recursively initialized. Whole aggregate deletion
has a zero-store representation theorem, while mapping payloads retain their
physical frame conditions. Scalar deletion preserves packed neighbors.
Source guards distinguish proved element removal/layout relocation from an
unknown lifetime history invalidated by raw writes or world replacement,
which yields an inconclusive observation.

Explicit reference casts retain the resolved location, aliases and slice
descriptors. Bytes-to-fixed-bytes conversion reads the used prefix and pads on
the right, excluding dirty unused suffixes. Source calldata-tail checks retain
native signed/wrapping behavior and distinguish semantic reverts from raw views
outside the physical typed domain. ABI copy decoding has its own check order.
Host fuel/materialization exhaustion is separate from source panic/revert and
from EVM resource adequacy.

The memory relation constrains used packets; unused byte padding and unallocated
fringe remain unconstrained. Padding-insensitivity theorems admit native dirty
storage-copy packets as well as the canonical interpreter representative.
Applying these theorems to compiler-emitted packets requires establishing the
representation relation for their concrete bytes.

The implementation separates executable copying from its proof clients.
\src{Operations.lean} implements allocation, conversion,
aliasing and physical stores. Its \src{Laws.lean} companion proves local
operation equations, and \src{Frames.lean} proves world and invocation frames.
Memory representation is split into predicates and observations
(\src{Base.lean}), scalar and pointer-cell stores
(\src{Cells.lean}), and constructor correctness
(\src{Construction.lean}), all under the aggregate directory.
The \src{Copy.lean} and \src{Representation.lean} modules re-export these
interfaces. Executable operations and proof clients share one heap and copy
interpreter.

\section{Related work and comparison method}
We compare the developments in Table~\ref{tab:related} along four dimensions:
executable domain, effect-order and alias representation, proved consumer, and
validation oracle. The comparison uses primary papers and, where available,
selected interfaces from the cited implementation revisions. Evaluation figures
for other projects are taken from their authors' reports.

\begin{table}[t]
\caption{Positive claims established by inspected primary artifacts.}
\label{tab:related}
\small
\begin{tabularx}{\columnwidth}{@{}>{\raggedright\arraybackslash}p{.27\columnwidth}X@{}}
\toprule
Development & Relevant semantic or proof interface \\
\midrule
K Solidity~\cite{ksolidity} & Executable structural operational source semantics
in K; contract analysis clients. \\
Lolisa~\cite{lolisa} & Solidity subset in Coq with stronger typed syntax and
a reusable formal memory framework. \\
Isabelle/\allowbreak Solidity~\cite{isabelle,isabelle-safety,isabelle-shallow} & Deep executable semantics,
grammar fuzzing, program proofs and type safety; a separate shallow embedding with a
storage model and verification-condition generation. \\
Solc-verify~\cite{memory-smt} & SMT-oriented encoding of storage pointers,
memory references, allocation and copying. \\
SolidiKeY~\cite{memory-key,solidikey-lean} & Logical data-area formalization;
Lean fragment compilation into a word-addressed machine under explicit input
relations. Physical EVM layout and fork/gas correspondence are separate. \\
SolidCore~\cite{solidcore} & Lean source execution with captured storage
references and documented intrinsic compiler-specific child orders;
reference-toolchain replay. \\
FORYU~\cite{foryu} & Rocq Yul semantics with a proved liveness-analysis checker
and reported evaluation on more than 1,500 contracts. \\
Solidity Spec & Per-group child choices, physical packet/alias conditions,
checked resolved entries and separate source/model/resource observations;
explicitly scoped safety and target-proof obligations. \\
\bottomrule
\end{tabularx}
\end{table}

Hajdu and Jovanovi\'c encode local storage pointers and memory allocation/copying
for solver automation~\cite{memory-smt}; da Silva et al. study object identity
and proof-oriented data-area representations~\cite{memory-key}. The inspected
SolidiKeY Lean compiler theorem relates a source fragment to a word-addressed
machine under explicit certificate and input-relation premises~\cite{solidikey-lean}.
Our representation relations track byte-addressed heap packets and ordered
physical writes, connecting them to guarded Cancun traces in the pilots.
Thornton and Marmsoler prove type safety for Isabelle/Solidity, including
consistency of pointer structures~\cite{isabelle-safety}. In our continuation
machine, the safety invariant connects physical heap representation to captured
targets, saved continuations, and the invocation checkpoint across child choices.

SolidCore is especially close in language and tooling. Its inspected order
design fixes construct-specific choices for a compiler configuration. Our
machine exposes independent choices for each eager group; its safety theorem
must therefore preserve the invariant under every permitted choice in the
profile. Isabelle's conformance methodology and FORYU's
semantics-plus-checked-consumer evaluation
inform our evaluation design~\cite{isabelle,foryu}. The remaining research
question is whether our combined interfaces reduce the assumptions and proof
effort of compiler-preservation compositions. Evaluating complete compiler passes
would extend the evidence beyond the two bounded pilots presented here.

\section{Package and reproducibility}
The package contains \src{Contract/} and \src{Solidity/}, a Lake configuration,
and pinned toolchain and dependency files. The default libraries are
\code{Contract} and \code{SoliditySemantics}; the only library dependency is
Lean's bundled \code{Std}. The release checks also require Python 3's
standard library. The pinned compiler and Anvil are additional dependencies
for external conformance checks.

From a Solidity Spec checkout with Lean 4.29.0 installed:
\begin{verbatim}
make replay JOBS=4
make paper
\end{verbatim}
The first command builds the package, runs the audit's negative controls,
enforces the theorem manifest and definition-source binding, then freshly
rechecks the library and three pilot roots plus scalar regressions sequentially
with \code{leanchecker --fresh}.
The second command uses Tectonic to rebuild this paper. CI runs the build and
fresh replay and retains diagnostic logs. \src{RELEASE.md} documents these
commands, and \src{THEOREMS.md} explains the theorem manifest. For external checks, set
\code{SOLC} to the pinned compiler executable and \code{LEAN\_BIN} to the Lean binary directory:
\begin{verbatim}
python3 scripts/check_conformance.py \
  --solc "$SOLC" --lean-bin "$LEAN_BIN" \
  --output /tmp/new-conformance-run
\end{verbatim}
\src{CONFORMANCE.md} documents observation scope and the required tools.

The public runtime interface is \src{Solidity.Execution}. Its \src{Proofs}
and \src{Safety} submodules expose laws and admitted-configuration safety.
\src{Solidity.Compatibility.Execution} exports certificate correspondences.
These interfaces separate runtime operations, representation assumptions and
proof consumers.

\section{Limitations and future research}
This release is an incomplete formalization. Resolved declarations, source
translation and exact artifact binding are explicit assumptions. The safety
profile is narrower than the executable aggregate domain, and the target
interpreter covers a Cancun subset. The pilots relate prepared source entries
to manually paired bytecode fragments; verified compiler emission and general
lowering remain open. Compatibility testing covers selected IR cases, with
legacy elaboration left for future work. Generic terminal typing remains
existential, and the uniform host bound is far larger than practical executions
usually require.

Future research includes verified source and declaration elaboration, safety
for additional aggregate representations, linked external execution, and
broader ABI and assembly semantics. Compiler-preservation compositions require
an explicit observation and gas model and a justified connection to emitted
bytes. Further work includes a root-result-indexed terminal invariant,
state-dependent resource bounds, verified conversion placement, broader
independent execution coverage and proofs for emitted compiler artifacts.

\appendices
\section{Artifact map}
Short filenames in the text identify declarations in their semantic
layer. Table~\ref{tab:layers} locates those layers. The theorem manifest records
fully qualified names and defining modules. \src{SAFETY.md}, \src{PILOT.md},
\src{THEOREMS.md} and \src{CONFORMANCE.md} index the precise proof profiles,
artifact bindings, audit commands and external checks.

\begin{table}[t]
\caption{Representability and current certificate coverage are distinct.}
\label{tab:layers}
\small
\begin{tabularx}{\columnwidth}{@{}p{.25\columnwidth}X@{}}
\toprule
Layer & Executable meaning and boundary \\
\midrule
Typed contract & Indexed variables and slots; finite helper bodies and logs.
Parsing, ABI and gas are handled externally. \\
Certificate source & Resolved word expressions; restricted methods, ABI guards,
word/Boolean storage, branches and initialization. Closed parser admission is trusted. \\
Resolved machine & Physical aggregates, effects, loops and internal recursion.
Resolved typed input; safety covers the stated structural profile. \\
Yul expressions & Word operators and environments used by optimizer proofs;
restricted to the expression fragment. \\
Interaction & Memory/world operations, call oracle, creation executor and
fuel, with conditional reasoning about callees. \\
Bytecode & Partial raw-byte Cancun semantics. Unsupported instructions and
exhausted interpreter fuel are explicit non-success results. \\
\bottomrule
\end{tabularx}
\end{table}

\begin{thebibliography}{99}
\bibitem{artifact} OKcontract, ``Solidity Spec,'' standalone project.
\url{https://github.com/okcontract/solidity-spec}.
\bibitem{ksolidity} J. Jiao et al., ``Executable Operational Semantics of
Solidity,'' 2018. \url{https://arxiv.org/abs/1804.01295v1}.
\bibitem{lolisa} Z. Yang and H. Lei, ``Lolisa: Formal syntax and semantics for a
subset of the Solidity programming language in Mathematical Tool Coq,''
version 3, 2021. \url{https://arxiv.org/abs/1803.09885v3}.
\bibitem{isabelle} D. Marmsoler and A. D. Brucker, ``Isabelle/Solidity: A deep
embedding of Solidity in Isabelle/HOL,'' \emph{Formal Aspects of Computing},
37(2), 2025. \url{https://doi.org/10.1145/3700601}.
\bibitem{isabelle-safety} B. Thornton and D. Marmsoler,
``Type Safety for Isabelle/Solidity,'' author manuscript.
\url{https://marmsoler.com/assets/pdf/thornton_24_typesafety.pdf}.
\bibitem{isabelle-shallow} D. Marmsoler, A. Ahmed and A. D. Brucker,
``Isabelle/Solidity---A shallow Embedding of Solidity in Isabelle/HOL,''
\emph{Archive of Formal Proofs}, 2026.
\url{https://isa-afp.org/entries/Isabelle-Solidity.html}.
\bibitem{memory-smt} \'A. Hajdu and D. Jovanovi\'c, ``SMT-Friendly Formalization
of the Solidity Memory Model,'' ESOP 2020, pp. 224--250.
\url{https://arxiv.org/abs/2001.03256v2}.
\bibitem{memory-key} G. H. A. da Silva, W. Ahrendt and R. Bubel,
``Axiomatisation of Solidity Memory and Storage,'' SEFM 2025, LNCS 16192,
pp. 73--90, 2026. \url{https://doi.org/10.1007/978-3-032-10444-1_5}.
\bibitem{solidikey-lean} SolidiKeY authors, ``Solidity Lean,'' revision
\src{23ba92a3efbc73dce04ad0483294d42757375a0d}, inspected 11 November 2026.
\url{https://github.com/SolidiKeY/solidity-lean/tree/23ba92a3efbc73dce04ad0483294d42757375a0d}.
\bibitem{solidcore} Paradigm authors, ``SolidCore,'' revision
\src{1f6fe8636d023926792faf4bd024bba943492a63}, inspected 11 November 2026.
\url{https://github.com/paradigmxyz/solidity-lean/tree/1f6fe8636d023926792faf4bd024bba943492a63}.
\bibitem{foryu} E. Albert, S. Genaim and E. Martin-Martin,
``Towards Formally Verified Smart Contracts Compilation,''
FM 2026, LNCS 16557, pp. 255--274.
\url{https://doi.org/10.1007/978-3-032-26220-2_13}.
\bibitem{control} Solidity authors, ``Expressions and Control Structures,''
Solidity 0.8.36 documentation.
\url{https://docs.soliditylang.org/en/v0.8.36/control-structures.html}.
\bibitem{ir-changes} Solidity authors, ``Solidity IR-based Codegen Changes,''
Solidity 0.8.36 documentation.
\url{https://docs.soliditylang.org/en/v0.8.36/ir-breaking-changes.html}.
\bibitem{abi} Solidity authors, ``Contract ABI Specification,'' Solidity 0.8.36.
\url{https://docs.soliditylang.org/en/v0.8.36/abi-spec.html}.
\bibitem{layout} Solidity authors, ``Layout of State Variables in Storage and
Transient Storage,'' Solidity 0.8.36.
\url{https://docs.soliditylang.org/en/v0.8.36/internals/layout_in_storage.html}.
\bibitem{types} Solidity authors, ``Types,'' Solidity 0.8.36.
\url{https://docs.soliditylang.org/en/v0.8.36/types.html}.
\bibitem{eip2200} W. Tang, ``EIP-2200: Structured Definitions for Net Gas Metering.''
\url{https://eips.ethereum.org/EIPS/eip-2200}.
\bibitem{eip2929} V. Buterin and M. Swende, ``EIP-2929: Gas cost increases for
state access opcodes.'' \url{https://eips.ethereum.org/EIPS/eip-2929}.
\bibitem{eip3529} V. Buterin and M. Swende, ``EIP-3529: Reduction in refunds.''
\url{https://eips.ethereum.org/EIPS/eip-3529}.
\bibitem{eip170} V. Buterin, ``EIP-170: Contract code size limit.''
\url{https://eips.ethereum.org/EIPS/eip-170}.
\bibitem{eip3541} A. Beregszaszi et al., ``EIP-3541: Reject new contract code starting
with the 0xEF byte.'' \url{https://eips.ethereum.org/EIPS/eip-3541}.
\bibitem{eip3860} M. H. Swende et al., ``EIP-3860: Limit and meter initcode.''
\url{https://eips.ethereum.org/EIPS/eip-3860}.
\end{thebibliography}
\end{document}
